\chapter{Diseño e ingeniería de la solución fotovoltaica}
\label{ch:diseno-fv}

% BORRADOR 27/09/2026 (redactor-humanizador). Pendiente: revisión de Mario → humanizer → compilación.
% Estado de las decisiones: PROPUESTA pendiente de validación por Mario (18 pendientes abiertos en
% §8 de Memoria/analisis/seleccion-modulos-por-variante.md). Fuente única de cifras: ese documento.
% Falta en este capítulo:
%   - 6.1: subapartados de inversores y de estructura soporte (sin decidir; pendiente 18).
%   - 6.2 Dimensionado del campo generador: strings y ratio DC/AC.
%   - 6.3 Cálculos eléctricos y esquema unifilar.
%   - 6.4 Simulación en PVsyst: pérdidas detalladas.
%   - 6.5 Producción esperada: PR, P50 y P90.
%   - 6.6 Grado de autoconsumo y autosuficiencia.
%   - 6.7 Monitorización y verificación del rendimiento.
% TODO Mario: las referencias a capítulos y apartados todavía no escritos van en texto llano;
% convertirlas en \ref{} cuando existan.

\section{Selección de componentes: módulos, inversores y estructura}
\label{sec:diseno-componentes}

\todo[inline]{Borrador. Las elecciones de módulo de este apartado son una propuesta pendiente de validación y pueden cambiar con los resultados de la simulación y la comprobación estructural.}

En este apartado se seleccionan los componentes principales del generador fotovoltaico de cada una de las variantes definidas en el capítulo 5. Por el momento se desarrolla solo la selección del módulo, ya que las elecciones de inversor y de estructura soporte dependen en parte de ella y se tratarán más adelante en este mismo capítulo.

Antes de entrar en los productos concretos, se resume de forma breve el panorama de tecnologías de generación fotovoltaica disponible en el mercado, puesto que dicho panorama es el que explica qué candidatos se comparan en cada variante y por qué.

\subsection{Módulos fotovoltaicos}
\label{subsec:diseno-modulos}

\subsubsection{Tecnologías de generación fotovoltaica}
\label{subsubsec:diseno-tecnologias}

La práctica totalidad de los módulos comerciales para aplicaciones como la de este proyecto emplea células de silicio cristalino. \textbf{[pendiente: cuota de mercado del silicio cristalino, con fuente citable]} El silicio de grado solar se obtiene a partir de silicio de grado metalúrgico y, según el proceso de cristalización, da lugar a obleas policristalinas (método Siemens) o monocristalinas (método Czochralski) \cite{Perpinan2013}. Para formar el módulo, las células se encapsulan en EVA entre un vidrio frontal y una lámina posterior o un segundo vidrio, y el conjunto se cierra con un marco de aluminio anodizado \cite{Perpinan2013}.

La evolución de la célula en los últimos años se aprecia bien en la eficiencia. El manual de referencia situaba en 2013 la eficiencia de las células industriales de silicio entre el 13 y el 17~\% \cite{Perpinan2013}, mientras que todos los módulos candidatos de este trabajo declaran eficiencias de módulo de entre el 22,0 y el 24,8~\% en sus fichas técnicas. \textbf{[pendiente: dar de alta las fichas técnicas con bibliotecario-pdf para poder citarlas]} Dicha mejora procede sobre todo del paso de la oblea de tipo p a la de tipo n y de nuevas arquitecturas de célula. Entre los candidatos analizados aparecen tres:

\begin{itemize}
    \item \textbf{TOPCon}, célula de tipo n con contacto pasivado, que es la arquitectura de la mayoría de los módulos de gran formato comparados.
    \item \textbf{Heterounión (HJT)}, célula de tipo n que presenta el menor coeficiente de temperatura de los candidatos.
    \item \textbf{Contacto posterior}, con los contactos en la cara trasera y sin rejilla frontal, comercializada con las denominaciones ABC (AIKO) y HPBC (LONGi).
\end{itemize}

Como se puede observar en la tabla~\ref{tab:diseno-tecnologias-celula}, los coeficientes de temperatura de potencia de los candidatos se sitúan entre $-$0,24 y $-$0,30~\%/\textdegree{}C. Son valores claramente mejores que el $-$0,4~\%/\textdegree{}C que el mismo manual toma por defecto para el silicio cristalino cuando no se dispone de datos del fabricante \cite{Perpinan2013}. Se trata de un parámetro relevante en Madrid, ya que la temperatura de célula en verano reduce la potencia de todas las variantes.

\begin{table}[H]
    \centering
    \small
    \caption{Tecnologías de célula de los módulos candidatos. Elaboración propia a partir de las fichas técnicas de los fabricantes.}
    \label{tab:diseno-tecnologias-celula}
    \begin{tabular}{@{}p{3.3cm}p{5.2cm}p{1.8cm}p{2.2cm}@{}}
        \hline
        \textbf{Tecnología} & \textbf{Módulos candidatos} & \textbf{Eficiencia (\%)} & \textbf{Coef. temp. $P_{max}$ (\%/\textdegree{}C)} \\
        \hline
        Heterounión (HJT), tipo n & REC Alpha Pure-RX 470 & 22,6 & $-$0,24 \\
        Contacto posterior ABC, tipo n & AIKO Neostar 2P 475 (monovidrio) y 2P+ 475 (vidrio-vidrio) & 23,8 & $-$0,26 \\
        Contacto posterior HPBC & LONGi Hi-MO X6 465 y Hi-MO 9 670 & 22,8 y 24,8 & $-$0,28 y $-$0,26 \\
        TOPCon, tipo n & Trina Vertex S+ 440, Jinko Tiger Neo 525, JA Solar JAM66D45-630 y Trina Vertex N 620 & 22,0 a 23,61 & $-$0,30 a $-$0,29 \\
        c-Si semitransparente & Onyx Solar c-Si (35~\% de transparencia) y Polysolar PS-MC-ST-60 & -- & $-$0,32 y $-$0,36 \\
        \hline
    \end{tabular}
\end{table}

A la tecnología de célula se suman otras dos características que distinguen a los módulos actuales y que pesan de forma distinta según la variante:

\begin{itemize}
    \item \textbf{Composición del laminado.} El módulo monovidrio lleva un vidrio frontal de 3,2~mm y lámina posterior, mientras que el vidrio-vidrio sustituye dicha lámina por un segundo vidrio, de 1,6 a 2,0~mm en los candidatos comparados. En el vidrio-vidrio no hay lámina trasera que se degrade y el encapsulante retiene los fragmentos en caso de rotura. Por eso se exige en la variante de marquesinas.
    \item \textbf{Bifacialidad.} Con la cara trasera transparente, el módulo aprovecha también la radiación reflejada por el suelo. Los candidatos bifaciales declaran factores de bifacialidad del 75 al 80~\%. Sin embargo, la ganancia real depende del albedo y de que la cara trasera quede libre, por lo que solo se considera en las variantes donde se da esa condición.
\end{itemize}

La capa fina se descartó en el lucernario por densidad de potencia. Con una transparencia del 20~\%, la opción de silicio amorfo evaluada alcanza 34,0~Wp/m$^2$, lo que sobre la superficie disponible supondría en torno a 15,7~kWp. Dicha potencia no constituye una alternativa de dimensionado, sino una variante distinta.

\subsubsection{Método de selección}
\label{subsubsec:diseno-metodo-modulos}

En cada variante se comparan tres candidatos de una misma clase de potencia y formato, de fabricantes bien posicionados. Las características comparadas son las que pide la guía de la asignatura para justificar el módulo: potencia \ac{STC}, tecnología de célula, coeficiente de temperatura, degradación inducida por la luz (\ac{LID}), garantías, características especiales, posición en rankings y precio. A ellas se añaden los criterios propios de cada variante, que salen de lo que limita esa instalación en concreto: superficie, tensión de string, carga de viento, comportamiento frente al fuego o función arquitectónica del elemento.

Cada criterio se puntúa de 1 (peor) a 3 (mejor) y los empates se reparten a la media, siguiendo el método empleado en un trabajo previo del mismo máster. \textbf{[pendiente: decidir si el TFM de M. Arcas (2024) se cita; hoy figura como fuente no citable]} El precio se ordena por categorías de mercado de febrero de 2026 y no por modelo, de modo que sirve para comparar pero no para presupuestar. \textbf{[pendiente: dar de alta la fuente de precios pvXchange y la de bancabilidad PV-Tech Q4 2025]}

Cabe destacar que se ha optado por no emplear un módulo único para todo el proyecto. Con un módulo común, uno de los argumentos decisivos era la tensión de 1.500~V, que solo necesita la planta en suelo, y se estaría aplicando a instalaciones de cubierta y marquesina que trabajarán con inversores de 1.100~V. Elegir por separado permite que cada decisión se apoye en el criterio que limita su variante, a cambio de gestionar cuatro modelos de módulo en la simulación en lugar de dos. La tabla~\ref{tab:diseno-criterios-variante} resume dichos criterios.

\begin{table}[H]
    \centering
    \small
    \caption{Criterio limitante y candidatos comparados en cada variante. Elaboración propia.}
    \label{tab:diseno-criterios-variante}
    \begin{tabular}{@{}p{2.6cm}p{4.6cm}p{5.3cm}@{}}
        \hline
        \textbf{Variante} & \textbf{Qué limita la variante} & \textbf{Candidatos comparados} \\
        \hline
        V1a. Lucernario (\acs{BIPV}) & Función de lucernario: transparencia mínima del 20~\% & Onyx Solar c-Si 35~\%, Polysolar PS-MC-ST-60, AGC SunEwat Vision \\
        V1b. Cubierta plana & Superficie (1.582,5~m$^2$ netos de módulo) y entradas de continua del inversor & REC Alpha Pure-RX 470, AIKO Neostar 2P 475 monovidrio, LONGi Hi-MO X6 465 \\
        V2. Marquesinas & Huella (3.182,4~m$^2$ de módulo) y vidrio sobre personas y vehículos & AIKO Neostar 2P+ 475 vidrio-vidrio, Trina Vertex S+ 440, Jinko Tiger Neo 525 \\
        V3. Planta en suelo & Capacidad de acceso (en torno a 4,6~MW), no el suelo & LONGi Hi-MO 9 670, JA Solar JAM66D45-630, Trina Vertex N 620 \\
        \hline
    \end{tabular}
\end{table}

\subsubsection{Variante V1a: lucernario}
\label{subsubsec:diseno-modulo-v1a}

La justificación del lucernario no es energética sino arquitectónica: se busca mantener la luz natural sobre el hall y aportar control solar. Por ello se exige como requisito previo una transparencia publicada igual o superior al 20~\%.

Por puntos gana el Polysolar PS-MC-ST-60 (21,0 frente a 19,5 del Onyx), pero su única transparencia publicada es del 10~\%, por lo que no pasa el requisito. Esto puede deberse a que su alta densidad, de 199,3~Wp/m$^2$, procede precisamente de cubrir casi todo el vidrio con célula. Se propone en consecuencia el \textbf{Onyx Solar c-Si con un 35~\% de transparencia}, de 182~Wp y 1.650 × 850~mm, ya que es el único que publica una gama de transparencias y una transmitancia térmica, y el que mejor coeficiente de temperatura presenta de los tres. Con 330 módulos sobre 463,1~m$^2$ de vidrio activo se obtienen unos 60,1~kWp.

Como debilidades, dicho módulo tiene una garantía corta (5 años de producto y 80~\% de potencia a 10 años) y no publica el factor solar $g$ de su gama cristalina. Queda como reserva el Polysolar, condicionado a que confirme una versión con transparencia igual o superior al 20~\%.

\subsubsection{Variante V1b: cubierta plana}
\label{subsubsec:diseno-modulo-v1b}

La cubierta cuenta con una impermeabilización líquida continua, por lo que la estructura debe ser lastrada y no perforante. Dicha condición afecta a la estructura, no al módulo, al que solo pide que su formato encaje en los sistemas lastrados estándar a 10\textdegree{} y que su peso no desequilibre la sobrecarga. Ninguno de los candidatos lo hace: la diferencia de peso entre ellos, repartida sobre la cubierta, es de solo 0,2~kg/m$^2$.

Se propone el \textbf{AIKO Neostar 2P 475~W monovidrio}, que obtiene 27,5 puntos frente a 26 del REC y 25 del LONGi. Gana por poco, pero lo hace en los criterios que limitan la variante:

\begin{itemize}
    \item \textbf{Densidad.} Con una eficiencia del 23,8~\%, permite instalar 794 módulos y unos 377,1~kWp, frente a los 356,7~kWp del REC y los 360,4~kWp del LONGi en la misma superficie.
    \item \textbf{Longitud de string.} Con un inversor de 1.100~V admite 24 módulos por string, frente a los 14 del REC. Son 33 strings en lugar de 54, lo que reduce las entradas de continua, que fueron el límite de esta variante en el recálculo de superficies.
    \item \textbf{Degradación inicial.} Declara una \ac{LID} igual o inferior al 1~\%, frente al 2~\% del REC.
\end{itemize}

Queda como reserva el \textbf{REC Alpha Pure-RX 470}, módulo del trabajo previo sobre esta misma cubierta. Presenta tres ventajas nada menores en la cubierta de un \ac{CPD}: una carga trasera de ensayo de 4.000~Pa frente a 2.400~Pa, una clasificación europea de comportamiento al fuego exterior de cubierta BROOF(t2) que no declaran los otros dos, y un fichero de módulo ya probado en PVsyst. Si la comprobación de viento o la normativa de incendios aplicable a la cubierta exigen más de lo que declara AIKO, se pasa al REC y se acepta perder unos 20~kWp.

\textbf{[pendiente: los 794 módulos se estimaron con un factor de empaquetado validado con la geometría del REC; hay que rehacer la disposición de los sheds en la escena 3D con 1.757 × 1.134~mm]}

\subsubsection{Variante V2: marquesinas}
\label{subsubsec:diseno-modulo-v2}

En las marquesinas el módulo queda sobre personas y vehículos, por lo que se exige como requisito de paso que sea vidrio-vidrio. Se propone que la marquesina sea estanca mediante perfiles que actúen a la vez de fijación y de canalón, sin chapa bajo el módulo, de manera que la cara trasera quede libre y la bifacialidad tenga sentido.

Por puntos gana el Trina Vertex S+ 440 bifacial, con 27 frente a 25,5 del AIKO. Sin embargo, cuatro de los criterios que gana en solitario (precio, peso, corriente de cortocircuito y bancabilidad) son comerciales o de dimensionado del inversor, y ninguno toca lo que limita la variante. Se propone por ello el \textbf{AIKO Neostar 2P+ 475~W vidrio-vidrio}, por tres razones:

\begin{itemize}
    \item \textbf{Densidad.} Alcanza unos 758,6~kWp frente a 700,5~kWp del Trina en la misma huella, un 8,3~\% más. Para que el Trina igualase la energía anual, su ganancia bifacial simulada debería superar el 8~\% aproximadamente, y sobre asfalto no se dispone de un dato verificado que lo sostenga.
    \item \textbf{Seguridad sobre vehículos.} Es el único de los tres que declara clase A de fuego según IEC, y monta vidrio de 2 + 2~mm frente a 1,6 + 1,6~mm.
    \item \textbf{Robustez frente a la solución estructural.} Si la marquesina acaba llevando chapa bajo el módulo, la ventaja bifacial del Trina y del Jinko desaparece y la del AIKO se mantiene.
\end{itemize}

Cabe matizar que, con unos 758~kWp, las marquesinas vierten más excedente que con la potencia del recálculo de superficies, por lo que parte de dicha ventaja de densidad se perderá en el antivertido. \textbf{[pendiente: recalcular los excedentes de V2 con unos 758,6~kWp antes de usar la densidad como argumento cuantitativo]} Queda como reserva el \textbf{Trina Vertex S+ 440}, que pasaría a elegido si PVsyst da una ganancia bifacial superior al 8~\% sobre el pavimento real o si se prioriza la bancabilidad y el precio. Antes habría que confirmar su clase de fuego, que no figura en la ficha (no verificado).

\subsubsection{Variante V3: planta en suelo}
\label{subsubsec:diseno-modulo-v3}

En la planta de Colmenar Viejo el suelo sobra: la parcela tiene 58,84~ha y la planta ocuparía unas pocas hectáreas. Lo que limita es la capacidad de acceso, por lo que el criterio de densidad pesa poco y lo decisivo es la energía producida por kWp a lo largo de 30 años. Se fijan como requisitos de paso la tensión de 1.500~V y la bancabilidad AAA, que cumplen los tres candidatos \textbf{[pendiente: dar de alta la fuente PV-Tech Q4 2025]}, todos ellos bifaciales y vidrio-vidrio de 2.382 × 1.134~mm.

Se propone el \textbf{LONGi Hi-MO 9 LR8-66HYD-670M}, con 20 puntos frente a 17,5 del JA Solar y 16,5 del Trina. El orden se mantiene aun eliminando el criterio de densidad. Presenta el mejor coeficiente de temperatura ($-$0,26~\%/\textdegree{}C), la menor degradación anual (0,35~\%) y la mejor garantía de potencia (88,85~\% a 30 años). Su bifacialidad algo menor, del 75~\% frente al 80~\% de los otros dos, recorta parte de esa ventaja, y será la simulación la que determine cuánto. En cuanto a la característica antipolvo que declara el fabricante, no está cuantificada por terceros, por lo que no se ha tenido en cuenta.

Queda como reserva el \textbf{JA Solar JAM66D45-630/LB}, con bifacialidad del 80~\% y única ficha que declara la fijación para seguidor, de modo que pasaría por delante si la variante se estudia con seguidor a un eje. El Trina Vertex N 620, con la categoría de precio más baja, se deja como opción para el análisis de coste nivelado de la energía (\ac{LCOE}) del capítulo 9 si la diferencia de precio compensa la menor energía.

\subsubsection{Síntesis de la selección de módulos}
\label{subsubsec:diseno-modulos-sintesis}

La tabla~\ref{tab:diseno-modulos-variante} recoge la propuesta de módulo para cada variante. El detalle de las tablas de puntuación queda documentado en el trabajo de ingeniería del proyecto. \textbf{[pendiente: decidir si esas tablas completas pasan a un anexo]}

\begin{table}[H]
    \centering
    \small
    \caption{Propuesta de módulo fotovoltaico por variante. Elaboración propia.}
    \label{tab:diseno-modulos-variante}
    \begin{tabular}{@{}p{2.1cm}p{3.6cm}p{2.3cm}p{2.4cm}p{2.2cm}@{}}
        \hline
        \textbf{Variante} & \textbf{Módulo propuesto} & \textbf{Tecnología y potencia} & \textbf{Reserva} & \textbf{Potencia estimada} \\
        \hline
        V1a. Lucernario & Onyx Solar c-Si, 35~\% de transparencia & c-Si semitransparente, 182~Wp & Polysolar PS-MC-ST-60 (condicionada) & 60,1~kWp \\
        V1b. Cubierta plana & AIKO Neostar 2P 475~W monovidrio & ABC tipo n, 475~W & REC Alpha Pure-RX 470 & 377,1~kWp (provisional) \\
        V2. Marquesinas & AIKO Neostar 2P+ 475~W vidrio-vidrio & ABC tipo n, 475~W & Trina Vertex S+ 440 bifacial & 758,6~kWp \\
        V3. Planta en suelo & LONGi Hi-MO 9 LR8-66HYD-670M & HPBC bifacial, 670~W & JA Solar JAM66D45-630/LB & Limitada por el acceso \\
        \hline
    \end{tabular}
\end{table}

En V1b y V2 se repiten el fabricante y la célula, pero no el producto: monovidrio de 3,2~mm en la cubierta y vidrio-vidrio de 2 + 2~mm en la marquesina, cada uno elegido por criterios distintos. Ambas versiones comparten parámetros eléctricos, lo que permitiría emplear el mismo diseño de string y un único modelo de inversor de 1.100~V en las dos variantes. La planta en suelo trabajará a 1.500~V con otro inversor, y el lucernario, con un inversor propio de pequeña potencia.

Por otro lado, el REC Alpha Pure-RX 470 que se empleó en el trabajo previo sobre esta cubierta no se confirma en ninguna variante, aunque se mantiene como reserva de V1b. En V1b se descarta por densidad, con un 5,4~\% menos de potencia en la misma superficie, y por sus strings de 14 módulos, que multiplican las entradas de continua.

En conclusión, se propone un módulo distinto para cada variante, elegido según lo que limita cada instalación: la transparencia en el lucernario, la superficie y las entradas de continua en la cubierta, la seguridad del vidrio y la densidad en las marquesinas, y la energía por kWp en la planta en suelo. Las cuatro elecciones deben confirmarse con los ficheros de módulo en PVsyst, la comprobación de viento y la temperatura mínima del emplazamiento, que fijará la longitud definitiva de los strings en el apartado de dimensionado del campo generador.

\subsection{Inversores}
\label{subsec:diseno-inversores}

% TODO inversores: sin decidir. Se desarrollará tras validar los módulos. Partir de §6 de
% seleccion-modulos-por-variante.md (V1b y V2 a 1.100 V con string común; V3 a 1.500 V; V1a
% inversor propio de 50-60 kW). Ns calculados con Tmín = -10 °C PROVISIONAL (pendiente AEMET).
% La corriente de 30 A por MPPT del Huawei SUN2000-100KTL-M2 está marcada como NO VERIFICADA.

La elección del inversor de cada variante queda pendiente y se desarrolla en una versión posterior de este apartado.

\subsection{Estructura soporte}
\label{subsec:diseno-estructura}

% TODO estructura: sin decidir (pendiente 18 del documento de módulos y memoria
% project_estructura_soporte_pendiente). Debe cubrir las cuatro variantes: V1a perfilería
% existente del lucernario; V1b sistema lastrado no perforante compatible con la membrana líquida;
% V2 marquesina concreta (vía A o B, fabricante, inclinación, cimentación); V3 fija o seguidor
% a un eje. Decidir si este bloque sigue dentro del 6.1 o pasa a sección propia del capítulo 6.

La elección de la estructura soporte de cada variante queda pendiente y se desarrolla en una versión posterior de este capítulo.

% TODO: secciones 6.2 a 6.7 (ver cabecera del fichero).
